% page 619 of the facsimile (image p-618 of the scan)
% block 170.2 x 242.0 mm ; left margin 31.8 mm
\pgc{10.160mm}{0.000mm}{
\l{\cel{56}611}
\l{}
\l{\sou{vejetarar}. (\sou{netrans}.) Praktikar la nutro-sistemo (\textquotedbl{}da Pitagoro\textquotedbl{})}
\l{\cel{4}segun qua supresesas omna speci de karni o lia derivaji ne-}
\l{\cel{4}mediata. - DEFIRS.}
\l{}
\l{\sou{vekar}. (\sou{netrans.)}\cel{2}Cesar sua dormo. - DE.}
\l{}
\l{\sou{vekto.}\cel{2}Levero di la balanco, qua ocilas sur sua axo e suportas}
\l{\cel{4}la pleti. - L.}
\l{}
\l{\sou{vektoro}. (\sou{geom.}) Radio qua duktesas (en kurvo fokoza) de foko ad}
\l{\cel{4}irga punto di la kurvo. - DEFIRS.}
\l{}
\l{\sou{velo}. Tuko de stofo per qua la mulieri kovras sua vizajo por}
\l{\cel{4}shirmar su de la regardi, o por protektar su kontre atmosfero}
\l{\cel{4}kolda, kontre vento, e c. -\cel{2}EFIS.}
\l{}
\l{\sou{velino}. (\sou{tekn.})\cel{2}Felo de bovyuno, apretita quale pergameno, ma}
\l{\cel{4}plu delikata e plu glata, uzata por piktar miniaturi, por}
\l{\cel{4}skribar. - DEFIR.}
\l{}
\l{\sou{velito}. I. (\sou{che la Romani antiqua}). Infantriano kun armi lejera. -}
\l{\cel{4}II. (\sou{armei di Napoleon Iesma)}. Pedo-chasero. - EFIS.}
\l{}
\l{\sou{velkar}. (\sou{netrans.})\cel{2}Cesir esar fresha. (\sou{anke metaf.})\cel{2}D.}
\l{}
\l{\sou{velocipedo}. Sidilo sur roti, quan onu movigas per presar pedalo. -}
\l{\cel{4}DEFIS.}
\l{}
\l{\sou{velodromo}.\cel{2}Lico por la konkursi di biciklisti. - EF.}
\l{}
\l{\sou{veluro}. Stofo kun du strati horizontala de fili : la suba (reverso)}
\l{\cel{4}formacas texuro densa ; la sura (averso) esas sole dika (per}
\l{\cel{4}mikra slingi sekita ed igita interegale longa).- DEFIS.}
\l{}
\l{\sou{venar}. (\sou{netrans.}) (\sou{aludante la persono qua parolas}). Vidar (per-}
\l{\cel{4}sono, kozo) qua iras o portesas a la loko ube esas la parol-}
\l{\cel{4}ante). - FIS.}
\l{}
\l{\sou{vendar}. (\sou{trans., ad ulu)}. Cedar (ulo) (ad ulu) qua proprietos olu,}
\l{\cel{4}po pekunio-quanto (o po valorajo altra-sorta). - EFIS.}
\l{}
\l{\sou{veneno}. I. Liquido quan sekrecas ula animali per glando specala,}
\l{\cel{4}e qua, enduktite per morduro o pikuro aden la sango di altra}
\l{\cel{4}animalo, divenas nocanta o mortiganta. - II. (\sou{metaf.}) Dico}
\l{\cel{4}o doktrino danjeroza. - EFIS.}
\l{}
\l{\sou{venerala}.\cel{2}Qua koncernas la relati inter-sexua intima. - DEFIRS.}
\l{}
\l{\sou{veneracar}. (\sou{trans.})\cel{2}Traktar kun respekto profunda, religiala,}
\l{\cel{4}humila. - EFIS.}
\l{}
\l{\sou{venerdio.}\cel{1}La dio kinesma di la semano.}
\l{}
\l{\sou{veniala.(teologio katol.; pri peko)}. Pardonebla, pardoninda. -EFIS}
}
